\chapter{The component library: Abstract}
\label{s:components}
\index{Library!Components|textbf}

This chapter presents an abstract of the \MCX component library's
structure and categories. It intentionally does \emph{not} attempt to
enumerate every component with its parameter list -- the library is large
and changes between releases, so any hand-maintained table here would
inevitably drift out of date (as earlier editions of this manual did).
Instead, each component category chapter that follows (Sources, Optics,
Samples, Monitors, \ldots) documents a curated selection of representative
components in depth, with their full, current parameter lists pulled
directly from the component source's McDoc header at build time (see
section~\ref{s:mxdoc} for the McDoc comment format). For the complete,
authoritative, always-current list of every component and its parameters,
use the \verb+mxdoc+ front-end:
\begin{lstlisting}
mxdoc
mxdoc searchterm
mxdoc file.comp
\end{lstlisting}
The first form generates and opens a browsable \textit{index.html} catalog
of the whole library (both the standard installation and any local/custom
components); the second searches for and displays documentation for all
components matching \textit{searchterm}; the third displays the
documentation of one specific component file. See
section~\ref{s:mxdoc-run} in the User Manual for the full \verb+mxdoc+
reference.

\section{Component categories}
\label{s:comp-overview}

The library (located under the \verb+mcxtrace-comps+ directory of the
\MCX source/installation tree, see section~\ref{s:files} in the User
Manual for how \MCX locates it) is organised into the following top-level
categories:

\begin{description}
\item[sources] Components that define the initial photon ray state --
  including point, flat, divergent and full-featured laboratory-source
  models -- see chapter~\ref{c:sources}.
\item[optics] Mirrors, lenses, slits, crystal monochromators, and other
  beam-manipulating components.
\item[samples] Components modelling matter placed in the beam for
  scattering/absorption studies (SAXS models, powders, single crystals,
  \ldots) -- see chapter~\ref{c:samples}.
\item[monitors] Components that record/histogram the photon beam without
  otherwise affecting it.
\item[misc] Components that do not fit the other categories, e.g. the
  \texttt{Progress\_bar} and beam-statistics utilities.
\item[union] The Union framework: a set of process/geometry components
  that can be freely combined to describe complex, possibly nested and
  overlapping sample environments with multiple scattering, decoupling
  the description of \emph{where} matter is from \emph{what} it does to
  the photon (Compton/Rayleigh scattering, photoelectric absorption,
  powder diffraction, \ldots).
\item[sasmodels] Small-angle scattering form factors generated from the
  SasView/sasmodels project, exposed via the \texttt{SasView\_model}
  component.
\item[astrox] Components specific to the AstroX branch of \MCX, modelling
  grazing-incidence X-ray optics for astronomical instrumentation (nested
  Wolter-type mirror shells, pore and micropore optics).
\item[contrib] Components contributed by \MCX users. The \MCX core team
  provides the same technical support as for other categories, but
  authored and maintained the code less directly; contributed components
  are generally reliable, but see each component's own validation/bugs
  notes. Components in \texttt{contrib} that mature and gain wide use are
  periodically promoted into one of the core categories above -- always
  check with \verb+mxdoc+ rather than assuming a component's category from
  an older document.
\item[obsolete] Components that were renamed or superseded. They are kept,
  and continue to work, purely for backwards compatibility with existing
  instrument files; new instrument development should avoid them
  (\verb+mxdoc+ will point to the current replacement where one exists).
\end{description}

\section{Data files}
\index{Library!Components!data}

Some components require external data files, e.g. atomic form factors,
element-specific absorption/refraction data from the NIST database, or
multilayer/crystal reflectivity curves. Such files distributed with \MCX
are located in the \verb+data+ sub-directory of the library and are
accessible to any component using the shared library read routines
without needing local copies.

\begin{table}
  \begin{center}
    {\let\my=\\
    \begin{tabular}{|p{0.32\textwidth}|p{0.62\textwidth}|}
      \hline
       {\bfseries \MCX/data} & Description \\
       \hline
       Al.txt, He.txt, etc. & Data files from the NIST database (covering the elements with $Z\in[1..92]$) to be used for absorption and refraction.\\
       FormFactors.txt & Atomic form factors for elements with $Z\in[1..92]$.\\
       Ref\_W\_B4C.txt & Reflectivity of a W-B4C multilayer.\\
       Ref\_W\_Si.txt &  Reflectivity of a W-Si multilayer.\\
      \hline
    \end{tabular}
    \caption{Data files of the \MCX library.}
    \label{t:comp-data}
    \index{Library!Components!data}
    }
  \end{center}
\end{table}

\section{Component and instrument examples}

An overview of the example instrument library, organised by facility and
category, is given in the User Manual, chapter~\ref{s:instrument}. Every
\texttt{.instr} file there is discoverable via \verb+mxgui+'s
\textbf{New from Template} menu, or via \verb+mxdoc+.
